\documentclass[11pt]{article}

\usepackage[margin=1in]{geometry}
\usepackage{amsmath,amssymb,amsthm,mathtools}
\usepackage{enumitem}
\usepackage{booktabs}
\usepackage{tabularx}
\usepackage{microtype}
\usepackage[hidelinks]{hyperref}
\usepackage{xurl}

\newtheorem{definition}{Definition}
\newtheorem{lemma}{Lemma}
\newtheorem{remark}{Remark}

\title{Challenge Answer Audit Trails for Successor Maintenance:\\[0.5ex]
\large Leftward Reader Walks from Carrier Slots to Narrow Terminals}
\author{}
\date{Unified archive note v0.04 (2026-03-21; archive v482)}

\begin{document}
\maketitle

\begin{abstract}
Challenge answer carrier slots already solve the exact shipped location of one released audience/question answer.
But a later note that wants to audit that answer still has to spell out the leftward walk by hand: start at the exact carrier slot, recover the emitted answer object, recover the whole-question handoff, split it into one-piece capsules, follow each piece branch, and stop only at the approved narrow terminal owners.
This note gives that full walk one home.
It defines a compact \emph{challenge answer audit trail}: one maintained reader walk from one exact shipped answer carrier slot to the narrow terminal fields that discharge that audience/question answer.
\end{abstract}

\paragraph{Non-redundancy note.}
Synthesis~36 owns audience-keyed challenge routes, Synthesis~37 owns question-keyed stop profiles, Synthesis~38 owns piece-keyed challenge branches, Synthesis~41 owns exact piece-hook terminal witnesses, Synthesis~45 owns one-piece answer capsules, Synthesis~46 owns whole-question answer sheets, Synthesis~47 owns emitted answer cards, Synthesis~48 owns exact answer-carrier slots, Synthesis~50 owns the local smallest-sufficient answer-selection problem, Synthesis~54 owns the answer-side terminal citation normal form, Synthesis~74 owns the paired terminal citation discipline, Synthesis~75 owns the paired cutover rule, and Synthesis~17 owns the concrete worked instantiation.
This note owns only the assembled leftward audit walk once those narrower layers are already fixed and the live question really requires the full walk.

\paragraph{Scope.}
This is not a new route map, stop profile, answer capsule, answer sheet, answer card, carrier slot, or release stage.
It is a narrow audit-interface note about turning those already-owned objects into one audience/question reader walk.

\paragraph{Imports.}
Piece-keyed challenge branches are imported from Synthesis~38.\cite{series:challengebranches}
Exact piece-hook terminal witnesses are imported from Synthesis~41.\cite{series:challengeterminalwitnesses}
One-piece answer capsules are imported from Synthesis~45.\cite{series:challengeanswercapsules}
Whole-question answer sheets are imported from Synthesis~46.\cite{series:challengeanswersheets}
Emitted answer cards are imported from Synthesis~47.\cite{series:challengeanswercards}
Exact answer carrier slots are imported from Synthesis~48.\cite{series:challengeanswercarrierslots}
The concrete worked instantiation is imported from Synthesis~17.\cite{series:workedexample}
The smallest-sufficient answer-selection problem is imported from Synthesis~50, and the answer-side terminal citation normal form plus its paired reuse/cutover discipline are imported from Synthesis~54 and Synthesis~74--75 as the default later-paper stop rule around this full walk.\cite{series:challengeanswercitationdockets,series:challengeanswerreviewmenus,series:pairedterminalcitationdiscipline,series:pairedterminalcutoverrules}

\paragraph{Exports.}
This note exports (i) challenge answer audit trails, (ii) a carrier-to-terminal walk rule, (iii) a piece-completeness rule for whole-question audit walks, and (iv) a piece-local stop rule saying that later notes may cite one maintained leftward audit path only when the live issue is the exact full reader walk from one shipped answer location to the narrow terminal owners.
Otherwise later terse papers should default to the answer-side terminal citation normal form in Synthesis~54 together with the paired terminal discipline and paired cutover rule in Synthesis~74--75. Selecting whether a later note should cite this full walk or one of the narrower late-stage answer objects still belongs to Synthesis~50.\cite{series:challengeanswercitationdockets,series:challengeanswerreviewmenus,series:pairedterminalcitationdiscipline,series:pairedterminalcutoverrules}

\section{Why answer audit trails deserve their own note}
Knowing where a released answer lives is not yet the same thing as knowing how to audit it.
The carrier slot already tells a reader which exact shipped field or reusable clause carries the answer.
But a referee or auditor often wants one more maintained object: for this audience and this question, what is the full leftward walk back to the narrow terminal owners, piece by piece?
Without that object, later notes keep restating some variant of ``start at this clause, recover the answer card, recover the whole-question sheet, then split by piece and follow the right branch for each piece.''
That is not just stylistic repetition.
It is a maintained walk, so the series benefits from giving it one narrow home.

\begin{definition}[Challenge answer audit trail]
Fix one concrete base$\to$successor comparison, one audience key $a$, and one downstream question key $q$.
A \emph{challenge answer audit trail} is a tuple
\[
T(a,q)=(a,q,K,C,S,\Pi)
\]
containing one imported answer-carrier-slot key $K$, one imported answer-card key $C$, one imported answer-sheet key $S$, and an ordered family
\[
\Pi=(\pi_1,\dots,\pi_m)
\]
of piece traces, where the order of the piece traces agrees with the semantic-piece order declared by the imported answer sheet.
Each piece trace $\pi_i$ names exactly one imported answer capsule, one imported branch, the imported terminal-witness keys for that piece, and the approved narrow terminal field where that piece may stop.
\end{definition}

\begin{definition}[Carrier-to-terminal walk rule]
A challenge answer audit trail satisfies the \emph{carrier-to-terminal walk rule} if it starts at exactly one imported answer carrier slot and preserves that slot's exact carrier artifact and locator before moving leftward through the imported answer card, the imported answer sheet, the imported one-piece capsules, the imported piece branches, and the imported narrow terminal owners.
So the trail begins at the shipped answer location; it does not begin somewhere narrower and merely mention the carrier later.
\end{definition}

\begin{definition}[Piece-completeness rule]
A challenge answer audit trail satisfies the \emph{piece-completeness rule} if its ordered piece list agrees exactly with the piece order declared by the imported answer sheet and if it contains one and only one piece trace for each such semantic piece.
So a whole-question audit trail may not silently omit a support-only piece merely because that piece is not visible on the carried sentence surface.
\end{definition}

\begin{definition}[Piece-local stop rule]
A challenge answer audit trail satisfies the \emph{piece-local stop rule} if every piece trace names exactly one imported answer capsule and one imported branch for that audience/question piece, and if the branch endpoint is the approved narrow terminal field where audit stops for that piece.
The carrier slot is the beginning of the walk, not the endpoint of the proof.
\end{definition}

\begin{lemma}[Audit trails expose the full reader walk without changing ownership]
Assume every challenge answer audit trail satisfies the carrier-to-terminal walk rule, the piece-completeness rule, and the piece-local stop rule.
Then a later note may cite one audit trail as the full leftward reader walk for one audience/question answer without changing the ownership of carrier location, answer identity, whole-question handoff, one-piece assembly, or narrow terminal stop.
\end{lemma}

\begin{proof}
The carrier-to-terminal walk rule fixes the beginning of the walk at one exact shipped carrier slot.
The piece-completeness rule prevents a whole-question trail from becoming an incomplete or selectively surfaced sketch.
The piece-local stop rule then says that every semantic piece still terminates at the already-approved narrow field reached by the imported branch.
So the audit trail changes only packaging: it assembles one maintained reader walk from already-owned objects, but it does not re-own any of those component layers.\qedhere
\end{proof}

\begin{table}[t]
\centering
\small
\begin{tabularx}{\linewidth}{@{}l l X@{}}
\toprule
Layer & Canonical object & Question answered \\
\midrule
Exact carrier location & carrier slot & Which shipped field or reusable clause carries the released answer? \\
Full reader walk & audit trail & How does one reader walk leftward from that exact carrier slot to the narrow terminal owners, piece by piece? \\
\bottomrule
\end{tabularx}
\caption{Audit trails sit one level above carrier slots: they expose the full leftward reader walk, not a new answer object.}
\label{tab:answeraudittrails}
\end{table}

\section{Worked example pointer}
The maintained worked bundle now ships \path{example_successor_challenge_answer_audit_trails.json}.\cite{series:workedexample}
It records four audit trails for the canned Case-B successor.
The release-note public-status trail starts at the shipped wrapper field \path{example_successor_lineage_notice.json:status_summary}, then preserves the imported release-note answer card and answer sheet, and finally splits into a continuity-piece trace ending at \path{example_successor_continuity_verdict.json:continuity_decision} and a closure-piece trace ending at \path{example_publication_closure_verdict.json:closure_status}.
The referee public-status trail starts one level later at the reusable clause carrier but walks to the same narrow terminal owners.
The release-note and auditor compact-diff trails both start at the reusable compact-diff clause, then split into a visible numeric-delta piece, a visible changed-family piece, and a support-only recomputation piece that ends at \path{example_verifier_report.json:recomputed}.
Concrete keys already in the ledger include \path{release_note_public_status_sentence_answer_audit_trail} and \path{auditor_compact_diff_summary_answer_audit_trail}.

\paragraph{Why this helps the series.}
The late-stage stack is now cleaner at five distinct scales.
A capsule handles one semantic piece.
A sheet handles one whole audience/question handoff.
A card handles one released audience/question answer object.
A carrier slot handles where that released answer actually lives.
And an audit trail handles the full leftward reader walk from that exact location to the narrow terminal owners.
That keeps later papers short without letting the last mile of the audit path disappear into prose.

Most later terse papers should therefore stop at the answer-side terminal citation normal form and reopen this audit-trail note only when the exact full leftward reader walk is itself disputed.\cite{series:challengeanswerreviewmenus,series:pairedterminalcitationdiscipline,series:pairedterminalcutoverrules}

\paragraph{A minimal public answer-audit-trail card.}
\begin{table}[t]
\centering
\small
\begin{tabularx}{\linewidth}{@{}lX@{}}
\toprule
Field & Minimal public answer \\
\midrule
Release-note full-audit anchor & \nolinkurl{release_note_public_status_sentence_answer_audit_trail} keeps carrier slot \nolinkurl{release_note_public_status_sentence_answer_carrier_slot}, answer card \nolinkurl{release_note_public_status_sentence_answer_card}, answer sheet \nolinkurl{release_note_public_status_sentence_answer_sheet}, ordered pieces \nolinkurl{continuity_piece}/\nolinkurl{closure_piece}, and the two piece traces ending at \nolinkurl{example_successor_continuity_verdict.json:continuity_decision} and \nolinkurl{example_publication_closure_verdict.json:closure_status} \\
Auditor full-audit anchor & \nolinkurl{auditor_compact_diff_summary_answer_audit_trail} keeps carrier slot \nolinkurl{auditor_compact_diff_summary_answer_carrier_slot}, answer card \nolinkurl{auditor_compact_diff_summary_answer_card}, answer sheet \nolinkurl{auditor_compact_diff_summary_answer_sheet}, ordered pieces \nolinkurl{diff_delta_piece}/\nolinkurl{recompute_piece}/\nolinkurl{changed_family_piece}, and the three piece traces ending at \nolinkurl{example_compare_report.json:observed_diffs}, \nolinkurl{example_verifier_report.json:recomputed}, and \nolinkurl{example_successor_continuity_verdict.json:changed_line_items} \\
Selection rule & Stop here only while the live issue is the full maintained leftward walk from one exact carrier slot through card and sheet to every piece-local capsule / branch / narrow terminal owner, not merely the exact location or emitted-answer identity alone \\
Left/right boundary & Reopen Synthesis~48 when the live issue is still only exact shipped location; widen leftward to Synthesis~45--46 or piece-local branch owners only when the real dispute is a narrower imported capsule / sheet / branch rather than the maintained walk as a whole \\
Paired-terminal boundary & Once the full walk is no longer the moving part, later terse papers should usually stop at Synthesis~54 together with Synthesis~74--75 rather than reopen this note merely to restate a stable answer-side handoff or paired cutover judgment \\
\bottomrule
\end{tabularx}
\caption{A minimal public answer-audit-trail card. This is the smallest answer-side public answer later terse papers can quote directly when the live issue is the full maintained leftward walk from an exact carrier slot to the narrow terminal owners.}
\label{tab:minimal-public-answer-audit-trail-card}
\end{table}

This card is intentionally non-decorative: it pins the actual worked audit-trail keys, the exact imported carrier/card/sheet anchors, the ordered piece traces they preserve, and the full-audit stop rule without silently re-owning exact carrier location, emitted-answer identity, narrower capsule/branch disputes, or later review-menu / paired-terminal policy.\cite{series:workedexample}

\paragraph{A forced public answer-audit-trail matrix.}
\begin{table}[t]
\centering
\small
\begin{tabularx}{\linewidth}{@{}lX@{}}
\toprule
Field & Forced public answer \\
\midrule
Release-note trail floor & \nolinkurl{release_note_public_status_sentence_answer_audit_trail} still forces the same full walk only while the same carrier slot \nolinkurl{release_note_public_status_sentence_answer_carrier_slot}, answer card \nolinkurl{release_note_public_status_sentence_answer_card}, answer sheet \nolinkurl{release_note_public_status_sentence_answer_sheet}, ordered piece list, and piece traces still end at \nolinkurl{example_successor_continuity_verdict.json:continuity_decision} and \nolinkurl{example_publication_closure_verdict.json:closure_status} \\
Auditor trail floor & \nolinkurl{auditor_compact_diff_summary_answer_audit_trail} still forces the same full walk only while the same carrier slot \nolinkurl{auditor_compact_diff_summary_answer_carrier_slot}, answer card \nolinkurl{auditor_compact_diff_summary_answer_card}, answer sheet \nolinkurl{auditor_compact_diff_summary_answer_sheet}, ordered piece list, and piece traces still end at \nolinkurl{example_compare_report.json:observed_diffs}, \nolinkurl{example_verifier_report.json:recomputed}, and \nolinkurl{example_successor_continuity_verdict.json:changed_line_items} \\
Imported-owner floor & Those shortcuts remain honest only while every piece trace keeps the same capsule key, branch key, and approved terminal stop field for the same audience/question pair; the audit trail does not get to keep the same full-walk citation after any imported piece-level owner has changed \\
Completeness floor & The full-walk shortcut still requires the same ordered piece list and one piece trace per semantic piece; the matrix has already failed closed if a surfaced or support-only piece disappears, reorders, or acquires a different terminal endpoint \\
Staleness trigger floor & This matrix goes stale as soon as the carrier slot changes, the answer card changes, the answer sheet changes, a piece key reorders, a piece trace rebinds to a different capsule or branch, or any approved terminal stop field stops pointing at the same narrow owner \\
Escalation pointer & Stop at Synthesis~48 when the live issue is still exact location; stop here only while one maintained trail still forces the same whole walk; widen to Synthesis~45--46 or to the piece-local branch / terminal owners only after the real dispute has become the narrower imported object rather than the maintained walk itself \\
\bottomrule
\end{tabularx}
\caption{A forced public answer-audit-trail matrix. This is the smallest answer-side public answer later terse papers can quote directly when the live issue is whether one maintained audit trail still forces the same full leftward walk rather than merely which audit trail exists in the abstract.}
\label{tab:forced-public-answer-audit-trail-matrix}
\end{table}

This matrix is intentionally non-decorative: it pins the actual worked audit-trail keys that still force one exact carrier-to-terminal walk, together with the precise conditions under which that forcing remains honest or goes stale, without silently re-owning exact carrier location, emitted-answer identity, narrower piece-level ownership, or later review-menu / paired-terminal layers.\cite{series:workedexample}

\paragraph{One tiny answer-audit-trail drill.}
For the worked auditor compact-diff answer, \nolinkurl{auditor_compact_diff_summary_answer_audit_trail} still forces the three-piece walk from \nolinkurl{auditor_compact_diff_summary_answer_carrier_slot} through the shared answer card and sheet to the diff-delta, recomputation, and changed-family terminal stops at \nolinkurl{example_compare_report.json:observed_diffs}, \nolinkurl{example_verifier_report.json:recomputed}, and \nolinkurl{example_successor_continuity_verdict.json:changed_line_items}. If the live question narrows to ``where exactly is the shipped clause?'' the honest stop is no longer this note but Synthesis~48. If the live question stays on the full walk but one piece trace rebinds to a different capsule, branch, or terminal field, the matrix has gone stale and a later terse paper must reopen the exact audit-trail owner rather than keep quoting the old walk. Only once the full walk is fixed and the remaining issue becomes the answer-side maintained handoff or paired terminal cutover should the paper leave this note for Synthesis~54 or Synthesis~74--75.\cite{series:workedexample}

\begin{thebibliography}{99}\small

\bibitem{series:challengebranches}
Anonymous.
\newblock Challenge Branch Maps for Successor Maintenance: Piece-Keyed Escalation from Audience Routes to Approved Terminal Fields.
\newblock Series synthesis note (Synthesis~38), 2026. (This archive.)

\bibitem{series:challengeterminalwitnesses}
Anonymous.
\newblock Challenge Terminal Witness Ledgers for Successor Maintenance: Exact Piece-Hook Terminal Ownership for Certified Question Paths.
\newblock Series synthesis note (Synthesis~41), 2026. (This archive.)

\bibitem{series:challengeanswercapsules}
Anonymous.
\newblock Challenge Answer Capsules for Successor Maintenance: Audience-Question Minimal Answer Bundles from Certified Piece Hooks.
\newblock Series synthesis note (Synthesis~45), 2026. (This archive.)

\bibitem{series:challengeanswersheets}
Anonymous.
\newblock Challenge Answer Sheets for Successor Maintenance: Audience-Question Piece-Complete Handoffs from Answer Capsules.
\newblock Series synthesis note (Synthesis~46), 2026. (This archive.)

\bibitem{series:challengeanswercards}
Anonymous.
\newblock Challenge Answer Cards for Successor Maintenance: Audience-Question Released Answers from Answer Sheets and Sentence Locks.
\newblock Series synthesis note (Synthesis~47), 2026. (This archive.)

\bibitem{series:challengeanswercarrierslots}
Anonymous.

\newblock Challenge Answer Carrier Slots for Successor Maintenance: Exact Shipped Fields and Reusable Clauses from Answer Cards.

\newblock Series synthesis note (Synthesis~48), 2026. (This archive.)

\bibitem{series:challengeanswercitationdockets}
Anonymous.

\newblock Challenge Answer Citation Dockets for Successor Maintenance: Smallest-Sufficient Audience-Question Handoffs across Capsules, Sheets, Cards, Carrier Slots, and Audit Trails.

\newblock Series synthesis note (Synthesis~50), 2026. (This archive.)

\bibitem{series:challengeanswerreviewmenus}
Anonymous.

\newblock Challenge Answer Review Menus for Successor Maintenance: Smallest-Sufficient Referee Handoffs from Dockets, Profiles, Packets, and Ladders.

\newblock Series synthesis note (Synthesis~54), 2026. (This archive.)

\bibitem{series:pairedterminalcitationdiscipline}
Anonymous.

\newblock Paired Terminal Citation Discipline for Review Spines: Stable Checked-Answer Stops and Adjacent Request-Side Capstones.

\newblock Series synthesis note (Synthesis~74), 2026. (This archive.)

\bibitem{series:pairedterminalcutoverrules}
Anonymous.

\newblock Paired Terminal Cutover Rules for Review Spines: When to Stop, Widen, or Fail Closed.

\newblock Series synthesis note (Synthesis~75), 2026. (This archive.)

\bibitem{series:workedexample}
Anonymous.

\newblock Worked Example: Receipt Interlock for an Anonymous-DHT Reader-Privacy Claim.

\newblock Series synthesis note (Synthesis~17), 2026. (This archive.)

\end{thebibliography}
\end{document}
